\section*{Course Outline}

You should expect to spend approximately 8--20 hours per week reading your textbook, studying, and completing assignments for this class. Because this is an accelerated 8-week online course, students should plan to work consistently throughout the week and check Blackboard and KCTCS email daily. Details about the assignments are provided below.

\section*{Description of Assignments}

\begin{tcolorbox}[
  enhanced,
  colback=blue!2,
  colframe=DeepBlue,
  boxrule=0.8pt,
  arc=3pt,
  left=10pt,
  right=10pt,
  top=8pt,
  bottom=8pt
]
{\bfseries\color{DeepBlue} Course Point Summary}

\vspace{0.4em}

\renewcommand{\arraystretch}{1.22}
\begin{tabularx}{\linewidth}{@{}X c c r@{}}
\textbf{Assignment Category} &
\textbf{Number of Assignments} &
\textbf{Point Value} &
\textbf{Category Point Total} \\
\hline
Start Here Assignments & 5 & Varies & 37 \\
Interactive SmartBook Exercises & 12 & 20 & 240 \\
Chapter Practice Quizzes & 12 & 20 & 240 \\
Case Studies & 4 & 20 & 80 \\
Lab Exercises & 12 & 12 & 144 \\
Lab Quizzes & 12 & 12 & 144 \\
Module Exams & 4 & 100 & 400 \\
Final Review Quizzes & 8 & Varies & 20 \\
Final Exam & 1 & 100 & 100 \\
\hline
\textbf{Total} & \textbf{66} & -- & \textbf{1405} \\
\end{tabularx}
\end{tcolorbox}

\section{Start Here Assignments}

\textbf{These assignments are used as evidence of attendance during the first week of class. Because this is an accelerated 8-week online course, students must complete the first Start Here assignments early in the semester to remain enrolled and to gain access to the rest of the course.} \textit{Students who do not complete the required first-week attendance assignments may be dropped from the course without notification as per home campus procedures.}

The \textbf{Syllabus Quiz}, \textbf{Online Learning Survey}, and \textbf{Discussion Board} must be completed by \textbf{Thursday, June 18, 2026, at 11:59 p.m. eastern time}. Additional Start Here assignments, including \textbf{Succeeding in the Online Classroom} and \textbf{Fundamentals of Student Success}, must be completed by \textbf{Saturday, June 20, 2026, at 11:59 p.m. eastern time}.



\subsection{Below you will find a brief description of each of the Start Here Assignments:}

\textbf{Online Learning Survey (5 pts)} -- This is a true/false survey to indicate your readiness to succeed in an online course. The more true answers you have, the more likely you are to succeed in an online class. If you select false for one or more questions, you should carefully consider the time, technology, and study habits required to succeed in an online course.

\textbf{Syllabus Quiz (5 pts)} -- This important quiz ensures that you are familiar with important class policies before beginning your work. It must be repeated until you earn 5 out of 5 points. You will not have access to course assignments until you earn a perfect score on this quiz.

\textbf{Discussion Board: Let's get to know each other! (5 pts)} -- This discussion board assignment allows us to get to know each other. Please use this as an opportunity to connect with other students if you wish to establish an online or in-person study group.

\textbf{Succeeding in the Online Classroom (12 pts)} -- This assignment is designed to help you get familiar with the online learning environment. You will be introduced to the idea of a Learning Management System. That is the program you use to access your class. For us, that program is Blackboard.

\textbf{Fundamentals of Student Success (10 pts)} -- This assignment is designed to help you make the most of your studies of human anatomy and physiology by implementing effective study skills, time management strategies, organizational skills, understanding and maximizing your learning style, and how to effectively engage with the material you are studying, your classmates, and instructor.

\clearpage

\section*{Weekly Assignments}

Three to five assignments are due every 2--3 days every single week of the semester, \textbf{even during exam weeks, holidays, and breaks.} It is recommended that you complete them in the following order:

\begin{itemize}
  \item Read each chapter (0 points)
  \item Interactive SmartBook Exercises
  \item Chapter Practice Quizzes
  \item Case Studies
  \item Lab Exercises
  \item Lab Quizzes
\end{itemize}

\subsection*{Interactive SmartBook Exercises: Read and Answer Questions (20 points each, 240 points total)}

This course covers a total of 12 chapters. Each chapter begins with an assignment known as an Interactive SmartBook Exercise to help guide you through the important concepts in each chapter. There are a total of 12 SmartBook exercises worth 20 points each for a total of 240 points. These exercises have no time limit and may be repeated as many times as you wish before the due date to earn a perfect score for each assignment. Each exercise is designed to take the average student approximately 1 hour to complete, so schedule your time accordingly. Waiting until the last minute will reduce the number of attempts you can make at each exercise.

\begin{tcolorbox}[coursetip]
{\bfseries Important Note About SmartBook}

SmartBook assignments begin with questions, but you are not expected to answer everything from memory. As you work, you may switch back and forth between answering questions and reading the electronic textbook. The program is interactive and may recommend specific readings or topics for review based on your responses. Students often find that moving between the questions and the textbook helps them learn the material more efficiently.
\end{tcolorbox}

\subsection*{Chapter Practice Quizzes (20 points each, 240 points total)}

After reading and completing the SmartBook exercise for each chapter, you must complete a quiz over the same content covered in the SmartBook exercise. These quizzes will test your ability to remember and apply what you covered in SmartBook exercises. You may use your textbook and notes, but a 30-minute time limit will be imposed. You may repeat the quiz, beginning with an entirely new set of questions, up to 3 times. Your best score will be used for grading purposes. You will complete a total of 12 Practice Quizzes, 1 per chapter, worth 20 points each for a total of 240 points.

\subsection*{Case Studies (20 points each, 80 points total)}

Four Case Studies are assigned, one at the end of each module. They are designed to help you apply concepts from each chapter in a real-life situation. Responses are submitted in Blackboard. Each of the four case study assignments is worth 20 points for a total of 80 points.

\subsection*{Lab Exercises (12 points each, 144 points total)}

The laboratory portion of the course is presented in an entirely online format using Anatomy \& Physiology Revealed (APR), and Virtual Labs. Each of the 12 lab activities will involve one or both of those programs; some labs have two parts. The lab exercises may be completed an unlimited number of times before the due date and only your highest score will be used for grading. After the due date you may revisit the labs for review only. Later attempts will not increase your grade. Importantly, these lab exercises also prepare you for the weekly lab quizzes. Questions on the lab quizzes are taken directly from these assignments. Each lab exercise is worth 12 points for a total of 144 points.

\subsection*{Lab Quizzes (12 points each, 144 points total)}

Upon completing each lab activity, students may proceed to a lab quiz designed to assess understanding of key concepts from the laboratory activity. Lab quizzes are limited to 30 minutes. These quizzes may be repeated up to 3 times, only the highest score will be used for grading. A total of 12 lab quizzes are assigned worth 12 points each for a total of 144 points.

\section*{Exams}

\subsection*{Module Exams (400 points total; 240 from lecture, 160 from lab)}

Module exams will be assigned at the end of each 3-chapter module. Module assignments are due before the exam window opens. Exams will open according to the schedule posted in Blackboard and will remain available during the posted exam window. Exams can be taken any time during the exam window, but they cannot be repeated for higher scores within that time.

Each module contains a separate exam for lecture and lab. The lecture portion of the exam will be taken from material covered in each chapter of the text that you have practiced through SmartBook Exercises, Chapter Practice Quizzes, and case studies. The lecture portion of each module exam will contain 30 questions worth 2 points each for a total of 60 points. You will have 60 minutes to complete the lecture portion of each module exam.

The laboratory portion of the exam will cover concepts, skills, and memorization information contained in each laboratory exercise. The laboratory portion of each module exam will contain 20 questions worth 2 points each for a total of 40 points. You will have 40 minutes to complete the laboratory portion of each module exam.

The combined value of each lecture and laboratory module exam is 100 points. There are four module exams, so a total of 400 possible points for exams throughout the semester.

\subsection*{Final Review Quizzes (20 points total)}

The week before your final exam opens, you will have the opportunity to review concepts from throughout the course by completing final review quizzes. These review quizzes are divided by module and by lecture/lab content. Together, they are worth a total of 20 points. Per the discretion of your instructor, you may also have bonus assignments to help you review for the final.

\subsection*{Final Exam (100 points total; 60 from lecture, 40 from lab)}

The final exam will be a comprehensive exam covering all 12 lecture chapters and all 12 lab activities from the course. The lecture and laboratory portions will be administered separately. Both portions will be available during the final exam window posted in Blackboard. Your answers must be submitted before the deadline to earn credit for the final exam.

The final exam lecture portion will consist of 60 questions from lecture material worth 1 point each. You will have 120 minutes to complete this portion of the exam. The final exam laboratory portion will contain 40 questions from laboratory material worth 1 point each. You will have 80 minutes to complete the laboratory portion of the exam.